\documentclass[conference]{IEEEtran}

\usepackage[hyphens]{url}
\usepackage{amsmath}

\title{Beyond the Empire: A Leakage-Aware, Reproducible Benchmark and\\
False-Positive Audit for Historical Polity Area Change}

\author{
\IEEEauthorblockN{Shailesh Mishra, [Teammate 2], [Teammate 3], [Teammate 4]}
\IEEEauthorblockA{Department of Computer Science and Engineering\\
PES University, Bengaluru, India\\
7.shaileshmishra@gmail.com}
}

\begin{document}
\maketitle

\begin{abstract}
Machine learning models of historical state collapse are easy to build and hard to
trust: small samples, autocorrelated observations, undocumented selection, and
leaky feature construction routinely produce scores that look better than they are.
We present \emph{Beyond the Empire}, a reproducible benchmark and error-audit
pipeline for predicting territorial area change of historical polities, built on
pinned releases of the Seshat Global History Databank (Equinox and Polaris~2026)
and the Cliopatria geometry dataset. Rather than claiming a predictive success,
the project treats evaluation itself as the contribution: exact-window labels with
censoring instead of fabricated endpoints, lineage-grouped and temporal holdout
splits with fold-local preprocessing, majority-baseline framing, fold-specific
out-of-fold SHAP explanations, and a pre-specified confident false-positive audit
whose cases are queued for sourced historical review. On the 98-row century-scale
Equinox baseline (6 contractions), no model beats a majority-class baseline
(grouped macro-F1: forest 0.472, logistic 0.452, majority 0.483). On an expanded
756-observation, 50-year, three-class area-change panel (59 declines in 30
conservative groups), the best exploratory model improves grouped macro-F1 over
majority (0.327 vs.\ 0.279) but not on a time-plus-group holdout (0.287 vs.\
0.285), with paired bootstrap intervals spanning zero. We further quantify how
fragile the labels themselves are: under hypothetical $\pm$10\% independent
endpoint-area error, only 34 of 59 decline labels are invariant, and none of six
directly coded CrisisDB collapse records overlaps an eligible panel window. The
result is a measurement-aware benchmark, an honest negative-to-weak predictive
result, and a concrete, reviewable audit apparatus - all reproducible end to end
from pinned, hash-verified sources.
\end{abstract}

\begin{IEEEkeywords}
cliodynamics, historical databases, leakage-safe evaluation, explainable AI,
benchmark, false-positive audit, Seshat, Cliopatria
\end{IEEEkeywords}

\section{Introduction}
Predicting the decline of states and empires from historical data is an
attractive idea, and publicly available structured history - the Seshat Global
History Databank and its companion Cliopatria geometry dataset - makes it look
straightforward. It is not. A naive pipeline can leak future information into
features, treat repeated reconstructions of the same polity as independent
samples, silently censor exactly the collapse events it claims to study, and
report accuracy on a test set that shares lineages with training. Scores from
such pipelines are artifacts of design choices, not evidence about history.

This project started from that suspicion. Our aim was to build the \emph{rigor
apparatus} first - leakage-safe construction, honest baselines, and a
pre-specified audit of the model's most confident mistakes - and only then ask
what, if anything, the models learn. The deliverable is therefore not a claim
that historical decline is predictable, but a reproducible benchmark with an
explicit measurement audit, plus an interactive dashboard that keeps the
limitations visible next to every score.

Our contributions are:
\begin{enumerate}
\item A leakage-aware pipeline for historical area-change prediction over pinned,
hash-verified releases of Equinox, Polaris~2026, and Cliopatria, with exact-window
labels and censoring instead of interpolation or fabricated endpoints
(Section~III).
\item An evaluation design that separates related histories (conservative
source-linked groups) and time (purged temporal holdout), reports
majority-baseline-relative macro-F1, and treats repeated folds as correlated
(Section~IV).
\item A pre-specified false-positive audit: out-of-fold scores $\geq 0.70$ with a
non-decline outcome form a persistent, source-anchored review queue, explained
with fold-specific linear SHAP values that reconstruct each displayed score
exactly (Section~V).
\item A measurement audit of the labels and features themselves: release-to-release
label agreement, hypothetical label-sensitivity stress tests, a terminal-window
selection audit, a CrisisDB coverage check, and a feature time-support analysis
showing that strict dating leaves only 144 of 756 rows with any governance
measurement (Section~VI).
\item Honest negative-to-weak predictive results with full artifacts: 18 automated
tests, pinned SHA-256 inputs, and dashboard exports (Streamlit and a Power BI
project source candidate) (Section~VII).
\end{enumerate}

\section{Data}
\label{sec:data}
\subsection{Equinox baseline panel}
The century-scale baseline uses the Equinox release of the Seshat Global History
Databank (pinned commit, SHA-256 recorded in \texttt{run\_metadata.json}). One row
is one polity at one century mark. Treating the source's \texttt{Terr} variable as
log10 territory, an exact +100-year endpoint is required; rows without one are
censored, never labeled stable. This yields 98 labeled rows from 76 observed
polities spanning 3800~BCE to 1800~CE, with only 6 contractions at the 0.7
area-ratio threshold. Features are seven complexity/scale variables (hierarchy,
population, capital population, administrative levels, money, information,
governance indicators) plus exact previous-century differences; features above
60\% missing are removed, and indicators, median imputation and scaling are fit
inside each training fold only.

\subsection{Expanded 50-year panel}
The original proposal's 50-year, three-class target is not feasible from Equinox
(the exact 50-year lookup returns no matched rows, and the pipeline refuses to
interpolate). We therefore built a separate panel from the official Polaris~2026
workbook and Cliopatria v0.2.1 geometry, pinned by commit and SHA-256. We accept
only direct SeshatID links with exactly one active POLITY geometry at both $t$ and
$t+50$; ambiguous or composite geometries are rejected, never summed. Requiring a
future endpoint under the same ID yields 756 candidate rows (605 stable, 92 rise,
59 decline) across 113 conservative connected groups; only 30 groups contain a
decline. Observation origins are bounded so that both the current year and its
+50-year target fall no later than 1900~CE.

Two properties of this panel matter more than any score. First, it selects
\emph{continuing, mapped} polities: termination and identity transitions are
censored, so the panel can exclude exactly the collapse events the original
proposal wanted to study. Second, 592 registry-grid windows have a future endpoint
beyond their polity's registered end (267 with a unique current geometry) and are
therefore invisible to this design.

\section{Method}
\subsection{Labels}
For horizon $H$, the label compares log-area endpoints:
\[
y = \begin{cases}
\text{decline} & a_{t+H}/a_t \le 0.7\\
\text{rise} & a_{t+H}/a_t \ge 1.3\\
\text{stable} & \text{otherwise}
\end{cases}
\]
with $H{=}100$ years for the Equinox baseline and $H{=}50$ for the expanded panel.
The thresholds are operational study choices, not validated historical event
definitions; Section~VI quantifies their sensitivity. A missing future endpoint is
censored, and a registry end date is never treated as evidence of collapse.

\subsection{Leakage controls}
All preprocessing (missingness filter, indicators, imputation, scaling) is fit
inside training folds. Delta features use only the exact previous window; a
regression test proves that changing future feature values cannot change any
current model input. (An earlier internal audit diagnosed the time shift backwards
- shifting a source row at $s$ to merge key $s{+}H$ retrieves the \emph{past}, not
the future; the corrected direction is enforced by tests.)

\subsection{Splits and models}
Grouped evaluation uses stratified five-fold splitting repeated over three seeds,
grouped by observed polity identity (Equinox) or by conservative groups built from
exact predecessor/successor registry names and shared map names (expanded panel).
These groups are leakage safeguards, not verified lineage. The temporal holdout
trains only on rows whose label window ends before a cutoff and additionally
removes all test groups from training. Models are deliberately simple and
untuned: a stable-majority baseline, balanced logistic regression, and a shallow
random forest (150 trees, depth 3, min leaf 4). The majority baseline is always
reported, because with 6-59 positive rows it is the model to beat.

\subsection{Explanations}
Explanations are actual fold-specific linear SHAP values for the held-out logistic
models, in log-odds. A test verifies that per-row SHAP values plus the fold base
value reconstruct the displayed out-of-fold score to $10^{-8}$. Explanations are
associational - they describe the model, not historical causation - and feature
families are named literally (there are no fabricated military or succession
measures in the feature set).

\section{Results}
\subsection{Century-scale baseline}
Table~\ref{tab:equinox} reports the Equinox evaluation. No model beats the
majority baseline on grouped macro-F1, and per-class average precision for
contractions is computable on only 11 of 15 folds because several folds contain no
positive. Fold results are highly unstable, and repeated folds are correlated, so
standard deviations are not confidence intervals.

\begin{table}[t]
\centering
\caption{Equinox baseline, binary contraction at 100 years (98 rows, 6
positives). Macro-F1: grouped mean$\pm$sd over 15 folds; single temporal
holdout.}
\label{tab:equinox}
\begin{tabular}{lcc}
\hline\hline
Model & Grouped & Temporal \\
\hline
Majority & $\mathbf{0.483 \pm 0.012}$ & 0.484 \\
Logistic & $0.452 \pm 0.112$ & 0.466 \\
Forest & $0.472 \pm 0.046$ & $\mathbf{0.537}$ \\
\hline\hline
\end{tabular}
\end{table}

\subsection{Expanded 50-year panel}
Table~\ref{tab:fifty} reports the area-history models (current area, prior area,
past area change, historical age). The forest improves grouped macro-F1 over
majority, but the improvement vanishes on the time-plus-group holdout
(0.287 vs.\ 0.285), and paired group-bootstrap intervals on pooled seed-0
out-of-fold predictions include zero for both models (forest $+0.030$,
$[-0.014, 0.084]$; logistic $+0.028$, $[-0.021, 0.094]$). These are conditional
exploratory summaries, not proof of effect or of no effect.

\begin{table}[t]
\centering
\caption{Expanded 50-year three-class panel (756 rows, 59 declines, 113
groups). Grouped: macro-F1 mean$\pm$sd over 15 folds. Holdout: time-plus-group
split. Area-history forest.}
\label{tab:fifty}
\begin{tabular}{lcc}
\hline\hline
Model & Grouped & Holdout \\
\hline
Majority & $0.279 \pm 0.078$ & 0.285 \\
Logistic & $0.290 \pm 0.096$ & 0.226 \\
Forest & $\mathbf{0.327 \pm 0.060}$ & $\mathbf{0.287}$ \\
\hline\hline
\end{tabular}
\end{table}

A separately frozen exploratory scenario adds governance/scale features under an
explicit retrospective assumption (a single undated, undisputed source record is
assigned across the polity's registered lifetime, with provenance marked). Under
that assumption the combined forest reaches grouped macro-F1 0.375 and holdout
0.312, and its pooled improvement over the area-only forest has a positive
conditional bootstrap interval ($+0.063$, $[0.019, 0.108]$). However, its decline
average precision on the historical holdout drops from 0.411 to 0.173, the
polity-lifetime timing assumption is not verified, and one of three prespecified
source-narrative spot checks surfaced a possible temporal-scope mismatch (an
Eastern Han record whose description begins ``Western Han''). We therefore report
this as conditional exploratory signal, not successful forecasting.

\section{The False-Positive Audit}
\label{sec:fp}
The project's central artifact is an error audit, not a score. The rule was fixed
before inspection: any row with out-of-fold score $\geq 0.70$ and observed outcome
other than decline enters a persistent review queue. Fourteen cases currently meet
the rule; every one is tagged \emph{pending}, with fields for source URL, review
note, and a genuine-survival vs.\ label-artifact verdict that nobody has filled
in, because no sourced historical review has happened yet.

This design deliberately inverts the usual framing. The interesting output of a
historical decline model is not where it is right but where it is confidently
wrong: each case is either a polity that genuinely resisted contraction (a
candidate ``what protected them'' case study) or an artifact of identity, geometry,
or coding (a finding about the data). Distinguishing the two requires per-case
source review, so the pipeline exports the queue with reproducible context instead
of asserting verdicts. A fixed 120-record, outcome-blinded review package with
stratified sampling and blank adjudication fields is prepared for two independent
reviewers; no reviewer has been contacted, and no further model fitting is planned
until that measurement evidence exists.

\section{Auditing the Measurements}
\label{sec:measurements}
\subsection{Label sensitivity}
Historical boundaries carry unquantified error. As a stress test (explicitly not
an error estimate), we perturb both endpoint areas by independent multiplicative
bounds and ask whether each label's full ratio range stays in its class. Under
$\pm$5\%, 42/59 decline labels are invariant; under $\pm$10\%, 34/59; under
$\pm$20\%, 25/59 - and \emph{zero} of 605 stable labels survive $\pm$20\%.
Shifting both lookups by the same calendar offset leaves 41/59 decline labels
invariant at $\pm$1 year, 20/59 at $\pm$10 years, and 3/59 at $\pm$25 years.
Roughly half the decline labels are artifacts of threshold and timing choices at
plausible stress levels; any future predictive claim must survive this audit.

\subsection{Release and cross-source checks}
Consecutive pinned Cliopatria releases (v0.2.0, v0.2.1) agree on all 754 shared
polity-year labels - a limited stability check, not independent replication, since
both releases share underlying sources. The Polaris workbook's CrisisDB sheet
contains 166 crisis annotations, six with directly coded collapse; none of the six
overlaps an eligible panel window, and famous cases (Cahokia, Tiwanaku, Wari) lack
direct geometry in the source. The panel therefore cannot validate collapse
sensitivity against these records, which is itself a benchmark-level finding about
coverage.

\subsection{Feature time support}
Many Polaris feature records carry no dates. Under a strict policy (one dated,
undisputed, non-uncertain record covering the observation time), only 144 of 756
rows have any of 20 governance/scale features - administrative level has 10 rows,
military level 7, professional soldiers 13. The official website displays undated
records against the polity's full lifetime; adopting that display assumption as a
separate, provenance-marked scenario raises coverage to 727 rows, but it is a
display convention, not independently dated measurement. Time-support policy, not
model choice, is the binding constraint on governance-feature modeling.

\section{Dashboard Artifacts}
Two dashboard surfaces expose the results with their caveats attached. A tested
Streamlit app provides a polity browser (held-out score history, observed
territory, fold-specific SHAP explanations), the false-positive queue, evaluation
tables, and the separate 50-year research summaries. A Power BI project source
candidate (PBIP/PBIR with TMDL semantic model) reproduces the polity explorer,
the audit queue, and a score-by-year scatter; its report JSON passes Microsoft's
published schema validation, but it has not yet been opened and rendered in Power
BI Desktop, and we do not claim it has. All scores are labeled uncalibrated; the
interfaces state that scores are not modern risk probabilities and that
explanations are not causes.

\section{Limitations}
This is an exploratory study with a small, selected sample. Labels are area
change, not verified collapse, survival, institutional quality, or welfare. The
continuing-ID design censors terminations. Groups are conservative safeguards,
not verified lineage. Sources are retrospective reconstructions, so past-only
features prevent simple leakage but not hindsight bias. No calibration,
independent historical adjudication, prospective validation, external-region
validation, or causal design has been done, and no modern-country scoring is
performed or implied. The expanded-panel holdout has been inspected during
development and is not untouched confirmatory evidence.

\section{Future Work}
The ranked path to stronger evidence is measurement, not tuning: (i) two-reviewer,
source-cited adjudication of the prepared blinded review package and of the 14
queued false positives, with disagreement tracking; (ii) a reviewed
predecessor/successor mapping to replace conservative grouping; (iii) confirming
undated-record semantics with the source maintainers before any governance claim;
(iv) only then, a newly frozen protocol with held-out regions/periods and
dependence-aware intervals. The original proposal's modern bridge (V-Dem/World
Bank domain-shift check and a resolved-case retrospective such as the USSR) is
intentionally deferred until a construct-valid bridge exists; we deliberately
exclude scoring any currently governing country.

\section{Conclusion}
Beyond the Empire shows that the honest contribution of a student-scale historical
prediction project can be the evaluation itself: pinned data, leakage controls,
majority-relative metrics, exact reconstructions of every displayed explanation, a
pre-specified audit of confident errors, and a quantitative account of how fragile
the labels are. The models, on present evidence, do not beat doing nothing fancy.
The benchmark, the audit queue, and the measurement findings are the results.

\section*{Availability}
Code, tests, protocols, and dashboards: \url{https://github.com/ShalikS-74/Vdui_mini_project}
(branch \texttt{completion/audit-pipeline-dashboard}). Sources are pinned by commit
with SHA-256 verification; the Equinox source is CC BY-NC-SA.

\begin{thebibliography}{9}
\bibitem{cliopatria}
Cliopatria: a global geospatial dataset of polities, 3400 BCE--2024 CE,
\emph{Scientific Data}, 2025. [Online]. Available:
\url{https://www.nature.com/articles/s41597-025-04516-9}
\bibitem{scheffer}
M. Scheffer \emph{et al.}, ``The vulnerability of aging states: A survival
analysis across premodern societies,'' \emph{PNAS}, 2023. [Online]. Available:
\url{https://www.pnas.org/doi/10.1073/pnas.2218834120}
\bibitem{bennett}
J. S. Bennett, ``Retrodicting the rise, spread, and fall of large-scale states in
the Old World,'' \emph{PLOS ONE}, 2022. [Online]. Available:
\url{https://journals.plos.org/plosone/article?id=10.1371/journal.pone.0261816}
\bibitem{turchin-scale}
P. Turchin \emph{et al.}, ``Scale and information-processing thresholds in
Holocene social evolution,'' \emph{Nature Communications}, 2020. [Online].
Available: \url{https://www.nature.com/articles/s41467-020-16035-9}
\bibitem{polaris}
Polaris 2025 Seshat data release report, Stanford GSB. [Online]. Available:
\url{https://www.gsb.stanford.edu/faculty-research/publications/polaris2025-seshat-data-release-expanded-dataset-cliodynamic-research}
\bibitem{equinox}
Equinox data release, Seshat Global History Databank. [Online]. Available:
\url{https://github.com/seshatdb/Equinox_Data}
\bibitem{codebook}
Seshat codebook. [Online]. Available: \url{https://seshat-db.com/codebook}
\bibitem{selection}
Selection and missingness in Seshat, eScholarship. [Online]. Available:
\url{https://escholarship.org/uc/item/4b11h9k0}
\bibitem{psychohistoryml}
PsychohistoryML public repository (self-reported, unaudited). [Online]. Available:
\url{https://github.com/TheApexWu/psychohistoryML}
\end{thebibliography}

\end{document}
